\section{Robustness}\label{sec:robust}

\paragraph{Second asset.} We repeat the full analysis for ETH/USDT (Appendix
Tables~\ref{tab:validation_eth}, \ref{tab:forecast_eth} and \ref{tab:regimes_eth}). The measurement
results are almost identical. Hourly BVC imbalance has a correlation of \ETHbvcCorr{} with observed
imbalance, the in-sample information ceiling of return-based classifiers is \ETHceilIS{} on hourly and
\ETHceilISm{} on one-minute bars, Garman-Klass volatility has a correlation of \ETHgkCorr{} with
realised volatility, and the Corwin-Schultz estimate has a median of \ETHcsMedian{} basis points
against a tick of \ETHtick. Section~\ref{sec:panel} extends the measurement results to twenty coins.

The forecasting results for ETH differ in one respect. The technical-indicator model is the only one
with a positive out-of-sample $R^2$ (\ETHrtwoTI\%, Clark-West $p=\ETHcwpTI$), and adding
microstructure variables makes forecasts significantly \emph{worse}: the Diebold-Mariano $p$-value
for ALL against TI is \ETHincALLTIDMp, with ALL less accurate. No microstructure or flow model beats
the zero forecast. The best trading rule is again one that uses no microstructure information, the
\ETHbestRule{} rule on \ETHbestModel{} with a gross Sharpe ratio of \ETHbestSR, and its deflated
Sharpe ratio of \ETHbestDSR{} is below conventional thresholds. The
reversal after aggressor flow found for BTC does not replicate for ETH at conventional significance
levels: the two-hour coefficient is negative in the full sample, under standard fees and in both
halves, but about half as large as for BTC, with $t$-statistics between $-0.7$ and $-1.7$, and it is
positive during the zero-fee window (Table~\ref{tab:response_splits}). In ETH, BVC imbalance has a
positive coefficient at eight to twelve hours ($t\approx2.3$, Appendix Table~\ref{tab:response_eth});
the bar-return coefficient is similar, so this reflects return continuation rather than information
in flow.

\paragraph{Stability of the flow response.} Table~\ref{tab:response_splits} re-estimates
regression~\eqref{eq:response} in sub-samples. For BTC the two-hour reversal is of similar size
under standard fees, during the zero-fee window and in both halves of the sample, and remains when
the bar return is held fixed. It is therefore not an artefact of the promotion or of one market
phase.

\begin{table}[t]
\centering\caption{Stability of the return response to observed flow}\label{tab:response_splits}
\resizebox{\textwidth}{!}{\begin{tabular}{lrrrrr}
\toprule
 & & \multicolumn{2}{c}{Observed imbalance} & \multicolumn{2}{c}{Observed imbalance $\mid z_t$}\\
\cmidrule(lr){3-4}\cmidrule(lr){5-6} Sample & Hours & $k=2$ & $k=4$ & $k=2$ & $k=4$\\
\midrule
\multicolumn{6}{l}{\emph{BTC/USDT}}\\
Full sample & 41,977 & \ensuremath{-}1.30 (\ensuremath{-}4.0) & \ensuremath{-}1.33 (\ensuremath{-}2.6) & \ensuremath{-}1.41 (\ensuremath{-}3.7) & \ensuremath{-}1.55 (\ensuremath{-}2.6)\\
Standard fees & 35,809 & \ensuremath{-}1.37 (\ensuremath{-}3.9) & \ensuremath{-}1.40 (\ensuremath{-}2.5) & \ensuremath{-}1.52 (\ensuremath{-}3.7) & \ensuremath{-}1.59 (\ensuremath{-}2.5)\\
Zero-fee window & 6,168 & \ensuremath{-}1.52 (\ensuremath{-}1.7) & \ensuremath{-}2.06 (\ensuremath{-}1.4) & \ensuremath{-}1.86 (\ensuremath{-}1.7) & \ensuremath{-}3.26 (\ensuremath{-}1.7)\\
Jan 2022 - Jun 2024 & 21,506 & \ensuremath{-}1.49 (\ensuremath{-}2.9) & \ensuremath{-}0.98 (\ensuremath{-}1.2) & \ensuremath{-}1.72 (\ensuremath{-}2.9) & \ensuremath{-}1.36 (\ensuremath{-}1.5)\\
Jul 2024 - Oct 2026 & 20,471 & \ensuremath{-}1.27 (\ensuremath{-}2.9) & \ensuremath{-}1.63 (\ensuremath{-}2.4) & \ensuremath{-}1.33 (\ensuremath{-}2.6) & \ensuremath{-}1.69 (\ensuremath{-}2.1)\\
\multicolumn{6}{l}{\emph{ETH/USDT}}\\
Full sample & 41,977 & \ensuremath{-}0.53 (\ensuremath{-}1.1) & 0.55 (0.8) & \ensuremath{-}0.79 (\ensuremath{-}1.5) & 0.21 (0.3)\\
Standard fees & 35,809 & \ensuremath{-}0.67 (\ensuremath{-}1.3) & 0.68 (0.9) & \ensuremath{-}0.88 (\ensuremath{-}1.6) & 0.30 (0.3)\\
Zero-fee window & 6,168 & 0.46 (0.4) & \ensuremath{-}0.70 (\ensuremath{-}0.4) & \ensuremath{-}0.07 (\ensuremath{-}0.1) & \ensuremath{-}0.96 (\ensuremath{-}0.5)\\
Jan 2022 - Jun 2024 & 21,506 & \ensuremath{-}0.44 (\ensuremath{-}0.7) & 0.10 (0.1) & \ensuremath{-}0.64 (\ensuremath{-}0.9) & 0.23 (0.2)\\
Jul 2024 - Oct 2026 & 20,471 & \ensuremath{-}0.63 (\ensuremath{-}0.9) & 0.87 (0.8) & \ensuremath{-}0.97 (\ensuremath{-}1.3) & 0.03 (0.0)\\
\bottomrule
\end{tabular}
}
\tabnote{Coefficients (basis points per standard deviation) and Newey-West $t$-statistics from
regression~\eqref{eq:response} at horizons of two and four hours, by sub-sample. The last two columns
include the standardised bar return $z_t$ as a control.}
\end{table}

\paragraph{Forecast horizon.} Table~\ref{tab:horizons} repeats the walk-forward for horizons of
1, 8 and 24 hours. The pattern matches ETH: the technical-indicator model has small positive
out-of-sample $R^2$ at one, eight and 24 hours, the largest being \BTChzTIeight\% at eight hours
(Clark-West $p=\BTChzTIeightp$), while the microstructure feature sets never improve on it.

\begin{table}[t]
\centering\caption{Forecast horizon, BTC/USDT}\label{tab:horizons}
\begin{tabular}{lrrrr}
\toprule
Horizon & EN: TI & EN: MS$_{\text{proxy}}$ & EN: MS & EN: ALL\\
\midrule
1h & 0.018 (0.01) & \ensuremath{-}0.025 (0.73) & 0.000 (0.06) & \ensuremath{-}0.004 (0.05)\\
4h & \ensuremath{-}0.055 (0.08) & \ensuremath{-}0.139 (0.73) & \ensuremath{-}0.006 (0.19) & \ensuremath{-}0.025 (0.09)\\
8h & 0.109 (0.00) & 0.010 (0.15) & 0.007 (0.27) & 0.034 (0.04)\\
24h & 0.072 (0.03) & \ensuremath{-}0.572 (0.78) & \ensuremath{-}0.841 (0.90) & \ensuremath{-}0.108 (0.32)\\
\bottomrule
\end{tabular}

\tabnote{Out-of-sample $R^2$ against the zero forecast (\%), Clark-West $p$-value against zero in
parentheses.}
\end{table}

\paragraph{Estimation window.} Replacing the expanding window by a six-month rolling window, which
adapts faster to changing conditions at the cost of noisier estimates, makes every model worse,
with out-of-sample $R^2$ of about $-0.4\%$ and no significant Clark-West statistic (Appendix
Table~\ref{tab:rolling}). Whatever predictability exists is not a short-lived, local relation that a
shorter window would capture.

\paragraph{Feature groups.} Dropping one feature group at a time from the full elastic net
(Appendix Table~\ref{tab:dropgroup}) leaves the out-of-sample $R^2$ between about $-0.17\%$ and
$+0.02\%$. The only positive value arises when the volatility group is removed, and it is too small
to matter; no single group is responsible for the negative results.

\paragraph{Conditional predictability.} Predictability could be concentrated in particular states,
for example when volatility is high or liquidity is low. Appendix Table~\ref{tab:conditional} splits
the out-of-sample period by causal volatility and illiquidity terciles, trading session, weekday and
fee regime. Out-of-sample $R^2$ is close to zero in every cell and no split shows consistent gains
across models.
